\section{Introduction}\label{sec:intro}

A functor on normed spaces extends fiberwise to $C^n$ vector bundles, with its natural transformations, when $C^n$
families of linear maps induce $C^n$ families (Theorem~\ref{thm:bundle-lifting}). We ask when the bifunctor
$(E,F)\mapsto\Alt^k(E;F)$ of alternating $k$-linear maps over a nontrivially normed field $K$ carries analytic bundles
to analytic bundles. For normed spaces $E$, $D$ and $F$ this concerns
$A^k_{E,D;F}\colon\cL(E,D)\to\cL(\Alt^k(D;F),\Alt^k(E;F))$, $A^k_{E,D;F}(u)(m)=m\circ(u,\dots,u)$. Let
$j_{E,F}\colon\Alt^k(E;F)\to\Mult^k(E;F)$ be the isometric inclusion. Composed with the induced isometric inclusion
$j_{E,F}\circ{-}\colon\cL(\Alt^k(D;F),\Alt^k(E;F))\to\cL(\Alt^k(D;F),\Mult^k(E;F))$, the map $A^k_{E,D;F}$ is the
diagonal of a bounded $k$-linear map (Proposition~\ref{prop:ambient-polynomial} with $g=\id_F$), and
Theorem~\ref{thm:coefficient-descent} decides when such a diagonal descends analytically through $j_{E,F}\circ{-}$. By Proposition~\ref{prop:homogeneous-lift}, $A^k_{E,D;F}$ is analytic at some
point if and only if it has a \emph{lift}: a bounded $k$-linear map $Q\colon\cL(E,D)^k\to\cL(\Alt^k(D;F),\Alt^k(E;F))$
with $Q(u,\dots,u)=A^k_{E,D;F}(u)$.

\subsection*{The splitting question}
Let $\cN^k_E$ be the kernel of the quotient map $\widehat\wedge^k_E$ from the completed projective tensor power of $E$
onto its completed projective exterior power (Section~\ref{sec:duality}):
\begin{equation}\label{intro:eq:source}
 0\longrightarrow\cN^k_E\longrightarrow\widehat\otimes^k_\pi E\xrightarrow{\ \widehat\wedge^k_E\ }
 \widehat\Lambda^k_\pi E\longrightarrow0.
\end{equation}
If $k!\ne0$ in $K$, antisymmetrization divided by $k!$ splits~\eqref{intro:eq:source} naturally in $E$ and solves every
lift problem (Propositions~\ref{can:prop:natural} and~\ref{prop:split-lift}(1)). If $k!=0$ in $K$, one $p$-cycle of
the coordinates of $K^k$, $p=\operatorname{char}K$, already forbids natural splittings (Theorem~\ref{can:thm:monomial},
Corollary~\ref{can:cor:p-cycle}). The question becomes when
the class $\ob_k(E,D)$ obtained from~\eqref{intro:eq:source} vanishes. It vanishes exactly when $(E,D)$ is \emph{good
for Banach values in degree $k$}: $A^k_{E,D;F}$ is analytic for every Banach space $F$; the pair is \emph{good in
degree $k$} if this holds for every normed space $F$ (Definition~\ref{fib:def:good}, Theorem~\ref{intro:thm:A}).

The class vanishes under a \emph{source splitting} $(\mathrm S^k_E)$, a bounded section of $\widehat\wedge^k_E$; a
\emph{target splitting} $(\mathrm R^k_D)$, a bounded retraction of the antisymmetrization $\alt^k_D$; or a
\emph{parameter splitting}, a bounded projection of $\widehat\otimes^k_\pi\cL(E,D)$ onto the closed span of the powers
$u^{\otimes k}$, which makes $\cL(E,D)$ \emph{polarizable in degree $k$}. Relative splittings, decompositions and
block assemblies combine them. The class can be nonzero: \emph{forced asymmetry} makes lifts asymmetric
(Section~\ref{sec:canonical}), and over a finite residue field \emph{Ramsey forcing} turns this asymmetry into a
triangular pattern of coefficients that norm estimates forbid (Section~\ref{sec:sharpness}).

\subsection*{Main results}


\begin{introtheorem}[descent and one obstruction class]\label{intro:thm:A}
Let $E$ and $D$ be normed spaces and $k\ge1$.
\begin{enumerate}
\item For every normed space $F$, $A^k_{E,D;F}$ is $C^\infty$ on $\cL(E,D)$; it is analytic at one point if and only if
it has a lift, and then it is a continuous polynomial.
\item The following are equivalent: \textup{(a)} $(E,D)$ is good for Banach values in degree $k$; \textup{(b)}
$u\mapsto\Lambda^k_\pi u$ has a bounded $k$-linear lift $\cL(E,D)^k\to\cL(\Lambda^k_\pi E,\widehat\Lambda^k_\pi D)$;
\textup{(c)} $\ob_k(E,D)=0$, where $\ob_k(E,D)$ is the pushout of the class of~\eqref{intro:eq:source} along the
restriction $\ev^k_{E,D}|_{\cN^k_E}\colon\cN^k_E\to\Mult^k_0(\cL(E,D);\widehat\Lambda^k_\pi D)$ of the evaluation map
$\ev^k_{E,D}(x_1\otimes\dots\otimes x_k)(u_1,\dots,u_k)=u_1x_1\wedge\dots\wedge u_kx_k$, with $\Mult^k_0$ the bounded
$k$-linear maps with zero diagonal; \textup{(d)} some bounded $k$-linear map
$\Psi\colon E^k\to\Mult^k_0(\cL(E,D);\widehat\Lambda^k_\pi D)$ satisfies
$\Psi(x_1,\dots,x_k)=\ev^k_{E,D}(x_1\otimes\dots\otimes x_k)$ whenever two of the $x_a$ coincide.
\item Let $(e_t)_{t\in\Sigma}$ be a linearly independent family in $E$ with dense span $V$, indexed by a totally ordered
set $\Sigma$; every normed space has one. The \emph{ordering lift} $Q_<(u)$, for $u\in\cL(E,D)^k$, is the linear map
$\Lambda^k_KV\to\widehat\Lambda^k_\pi D$ sending $e_{t_1}\wedge\dots\wedge e_{t_k}$, $t_1<\dots<t_k$, to
$u_1e_{t_1}\wedge\dots\wedge u_ke_{t_k}$. Composed with the canonical map $\Lambda^k_KV\to\Lambda^k_\pi E$, every lift
$Q$ as in \textup{(b)} is $Q_<+\Phi$ for a correction $\Phi$, $k$-linear in $u$ with $\Phi(v,\dots,v)=0$ for
$v\in\cL(E,D)$ (Theorem~\ref{ob:thm:ordering}). The pair $(E,D)$ is good for Banach values in degree $k$ if and only
if, for some such $\Phi$ and some $C$, $\norm{Q_<(u)y+\Phi(u)y}\le C\norm{u_1}\cdots\norm{u_k}\norm y_\pi$ for all $u$
and $y\in\Lambda^k_KV$, where $\norm y_\pi$ is the projective seminorm.
\end{enumerate}
\end{introtheorem}

Theorem~\ref{intro:thm:A} collects Lemma~\ref{lem:closed-smooth}, Propositions~\ref{prop:ambient-polynomial}
and~\ref{prop:homogeneous-lift}, Lemma~\ref{fib:lem:test}, Theorem~\ref{ob:thm:obstruction},
Corollary~\ref{ob:cor:adjacent} and Theorem~\ref{ob:thm:ordering}.

\begin{introtheorem}[canonical splittings exist exactly when $k!\ne0$]\label{intro:thm:B}
Let $k\ge2$.
\begin{enumerate}
\item If $k!\ne0$ in $K$, then $\alt^k_E/k!$, $\wedge^k_D/k!$ and the symmetrizer $(1/k!)\sum_{\sigma\in S_k}\sigma$ on
$\widehat\otimes^k_\pi P$ split the source, target and parameter sequences naturally in $E$, $D$ and $P$. Every pair is
good in degree $k$, and $\Alt^k$ is analytic on every hom space.
\item Suppose $k!=0$ in $K$. For every norm on $K^k$ with bounded coordinate functionals, no section of
$\wedge^k_{K^k}$, retraction of $\alt^k_{K^k}$ or projection onto the symmetric tensors commutes with the permutation
matrix of one $p$-cycle of the coordinates. If $E$ has $k$ linearly independent bounded functionals and
$\widehat\Lambda^k_\pi D\ne0$, no lift of $u\mapsto\Lambda^k_\pi u$ is symmetric, no section of $\widehat\wedge^k_E$
takes sign-isotypic values and no retraction of $\alt^k_D$ is sign-equivariant.
\item If $k!=0$ in $K$, then $A^k_{E,E;F}$ is analytic at no point for some Banach spaces $E$ and $F$, so $\Alt^k$ is
analytic on all hom spaces, or on those of Banach spaces, if and only if $k!\ne0$. If $K$ is moreover complete and not
spherically complete, $A^k_{E,D;K}$ is analytic at no point for some nonarchimedean Banach spaces $E$ and $D$.
\item If $K$ is complete and nonarchimedean, $\Alt^k(-;K)$ is analytic on all hom spaces of normed spaces if and only
if $k!\ne0$ in $K$ or $K$ is spherically complete.
\item Let a group act on a set $I$, and on $K^{(I)}$ by permuting the unit vectors. On the algebraic powers of
$K^{(I)}$, $\wedge^k$ has an equivariant section and $\alt^k$ an equivariant retraction if and only if the stabilizer of
each $k$-subset of $I$ induces on it a group of order nonzero in $K$; an equivariant projection onto the span of the
pure powers exists if and only if the stabilizer of each $k$-multiset on $I$ induces on its support a group of order
nonzero in $K$. The obstruction to an equivariant section is a family of Tate classes. One element inducing a
permutation of order $p$ on some $k$-subset forbids all three equivariant splittings, also on completed powers of normed
spaces with a compatible biorthogonal system of vectors and bounded functionals.
\end{enumerate}
\end{introtheorem}

Theorem~\ref{intro:thm:B} collects Propositions~\ref{can:prop:natural} and~\ref{prop:split-lift},
Theorem~\ref{can:thm:monomial}, Proposition~\ref{can:prop:asymmetry}, Theorem~\ref{thm:operator-obstructions},
Corollary~\ref{cor:classification}, Theorems~\ref{thm:scalar-classification} and~\ref{can:thm:permutation},
Corollary~\ref{can:cor:p-cycle} and Proposition~\ref{can:prop:tate}. Corollary~\ref{cor:bundle-cases} lists the bundle
settings in which fiberwise alternating maps form an analytic bifunctor; Proposition~\ref{prop:shear},
Corollary~\ref{cor:shear-bundle} and Proposition~\ref{prop:tangent-obstruction} give analytic atlases with nonanalytic
induced alternating or cotangent atlases, the latter over complete fields of characteristic $p>0$ that are not
spherically complete, in degrees $k\ge p$; over complete fields, ambient analytic forms keep pullback, wedge products
and exterior derivatives (Theorem~\ref{thm:ambient-forms}).

\begin{introtheorem}[how the class vanishes]\label{intro:thm:C}
Let $E$ and $D$ be normed spaces and $k\ge1$.
\begin{enumerate}
\item The pair $(E,D)$ is good for Banach values in degree $k$ if $(\mathrm S^k_E)$ or $(\mathrm R^k_D)$ holds, if
$\cL(E,D)$ is polarizable in degree $k$, if $(E,D)$ has a splitting decomposition of degree $k$, or if $(E,D)$ is an
$\ell^1$ block assembly of pairs with uniformly bounded lifts in all degrees $j\le k$. For decompositions into
$\mathrm S^{\le k}$-split and $\mathrm R^{\le k}$-split pieces, two pieces suffice.
\item If $E$ has a linearly independent family, indexed by a totally ordered set, with dense span $V$ such that the
strict triangle projection $\cP_<$ is bounded on $V\otimes_\pi V$, for instance one with a lower $r$-estimate for some
$r<2$, then $(\mathrm S^n_E)$ and $(\mathrm R^n_E)$ hold and $E$ is polarizable in degree $n$, for every $n\ge1$.
\item If $K$ is nonarchimedean, $E$ has a bounded linearly independent family indexed by a totally ordered set, with
dense span on which $E$ is $\ell^2$-dominated, and $D$ has an equivalent $r$-norm, $r\in(2,\infty)$, or an
equivalent nonarchimedean norm, then $(E,D)$ is good for Banach values in every degree.
\item Suppose $\operatorname{char}K=p>0$. The least degree, if any, in which $(E,D)$ is not good, or not good for
Banach values, is a power of $p$. If $E^*$ and $D^*$ are infinite-dimensional, the set of degrees in which $(E,D)$ is
good, or good for Banach values, is $\{1,\dots,p^h-1\}$ for some $h\ge1$ or $\{1,2,\dots\}$. For every normed space
$G$, each of $(\mathrm S^p_G)$, $(\mathrm R^p_G)$ and polarizability in degree $p$ implies the same condition in every
degree.
\item If $K$ is nonarchimedean and $E$ has a bounded linearly independent family $(e_t)_{t\in\Sigma}$, indexed by a
totally ordered set, with dense span $V$ and coordinate functionals $x\mapsto x(t)$ uniformly bounded on $V$, then the
class of~\eqref{intro:eq:source} pushes out to zero along every bounded linear map from $\cN^k_E$ to a Banach space with
an equivalent nonarchimedean norm. In particular $\ob_k(E,D)=0$ if $\widehat\Lambda^k_\pi D$ has such a norm.
\end{enumerate}
\end{introtheorem}

Theorem~\ref{intro:thm:C} collects Proposition~\ref{fib:prop:split}, Corollary~\ref{ob:cor:mechanisms},
Theorem~\ref{cmb:thm:decomposition}, Propositions~\ref{cmb:prop:normal-form} and~\ref{cmb:prop:blocks},
Theorem~\ref{fib:thm:lower}, Proposition~\ref{spl:prop:ordering}, Theorem~\ref{spl:thm:relative-ordering},
Corollaries~\ref{spl:cor:relative-ordering-na} and~\ref{fib:cor:degrees}, and Theorems~\ref{cmb:thm:initial-segments},
\ref{cmb:thm:column}, \ref{cmb:thm:polarizable-height} and~\ref{ob:thm:ultrametric}. A bounded retraction of $j_{E,F}$
gives analyticity for that value space $F$ only (Proposition~\ref{prop:split-lift}, Theorems~\ref{thm:spherical}
and~\ref{fib:thm:ultra}).

\begin{introtheorem}[forcing]\label{intro:thm:D}
Assume condition~\textup{(L)}: $\kappa$ is a finite field of characteristic $p$, $K$ is a subfield of $\kappa((X))$ on
which the absolute value $\abs a=\varrho^{\ord a}$, $0<\varrho<1$, is nontrivial, and $k\ge p$.
\begin{enumerate}
\item For $2\le r\le\infty$, the pair $(\ell^r(\N,K),\ell^r(\N,K))$, with $c_0(\N,K)$ for $r=\infty$, is not good for
Banach values in degree $k$.
\item If symmetric sequence spaces $E$ and $D$ form a pair good for Banach values in degree $k$, then for some $C$, the
\emph{product triangle bound} \textup{(PTB)}
\[
 \norm{\cP_<\bigl((ax)\otimes(by)\bigr)}_{D\otimes_\pi D}
 \le C\norm{\mu_a}_{E\to D}\norm x_E\norm{\mu_b}_{E\to D}\norm y_E
\]
holds for all finitely supported $a,x,b,y$, where $\mu_a$ is multiplication by $a$. If moreover
$\norm\xi_D\le C'\norm\xi_r$ for all finitely supported $\xi$ and some $r\in(2,\infty]$, then $E$ is $\ell^2$-dominated.
\item If $E$ has an equivalent nonarchimedean norm with $\dim_{\Kh}\widehat E=\infty$, then the pairs
$(E,c_0(\N,\Kh))$ and $(c_0(\N,K),E)$ are not good for Banach values in degree $k$.
\end{enumerate}
\end{introtheorem}

Theorem~\ref{intro:thm:D} collects Theorems~\ref{fib:thm:ell2} and~\ref{obs:thm:ptb},
Proposition~\ref{obs:prop:l2-domination}, and Theorem~\ref{fib:thm:local} with its proof, all proved in
Appendix~\ref{app:forcing}; Lemma~\ref{obs:lem:normalization} reduces condition~\textup{(L)} to fields dense in a
Laurent series field. Over every nontrivially normed field $K$ of characteristic~$2$, polynomial identities replace
Ramsey's theorem: for $E=c_0(\N,K)$, $\linf(\N,K)$ and $\ell^r(\N,K)$, $2<r<\infty$, the pair $(E,E)$ is not good for
Banach values in any degree $k\ge2$, and neither $(\mathrm S^k_E)$ nor $(\mathrm R^k_E)$ holds
(Theorem~\ref{obs:thm:char-two}). Over Laurent series fields over infinite fields and over densely valued fields, the
case $E=c_0(\N,K)$ for odd $p$ and the case $E=\ell^2(\N,K)$ for $p=2$ are open (Question~\ref{q:quest:residue}). In
characteristic~$2$, Proposition~\ref{fam2:prop:char-two} removes condition~\textup{(L)} from
Theorems~\ref{fam2:thm:staircase} and~\ref{fam2:thm:groups}(4), from Theorem~\ref{fam2:thm:tangent}(3) except its last
sentence, and from Corollary~\ref{fam2:cor:table}(2) for $E=c_0(\N,K)$ and for $E=\ell^{r_2}(\N,K)$, $r_2>2$.

\begin{introtheorem}[good pairs in classes]\label{intro:thm:E}
Let $\kappa$, $\varrho$ and $K$ be as in condition~\textup{(L)}, and let good mean good for Banach values.
\begin{enumerate}
\item Let $E$ be a symmetric sequence space in which $c_{00}$ is dense. The following are equivalent: $(E,E)$ is good
in some degree $k\ge p$; $(c_0,E)$ is good in some degree $k\ge p$; $\cP_<$ is bounded on $E\otimes_\pi E$;
$(\mathrm S^k_E)$ holds for some $k\ge p$; $(\mathrm R^k_E)$ holds for some $k\ge p$; $E$ is polarizable in some degree $k\ge p$; $(E,D)$ and $(D,E)$ are good in every
degree for every normed space $D$. If they hold, $E$ is $\ell^2$-dominated and $\varphi_E(m)^2\ge c\,m\ln m$ for some
$c>0$ and all $m\ge1$, with $\varphi_E(m)$ the norm of the indicator of $\{1,\dots,m\}$.
\item Let $E$ be as in \textup{(1)} and $D$ a symmetric sequence space with an equivalent $r$-norm, $r\in(2,\infty)$,
or an equivalent nonarchimedean norm. Then $(E,D)$ is good in some degree $k\ge p$ if and only if \textup{(PTB)}
holds, if and only if $E$ is $\ell^2$-dominated, if and only if $(E,D)$ is good in every degree.
\item For $1\le r_1,r_2\le\infty$, with $c_0$ in place of $\ell^\infty$, the pair $(\ell^{r_1},\ell^{r_2})$ is good in
some degree $k\ge p$ if and only if it is good in every degree, if and only if $\min(r_1,r_2)<2$ or $r_1=2<r_2$.
\item If $E$ has an equivalent nonarchimedean norm and $k\ge p$, then $(E,D)$ is good in degree $k$ for every normed
space $D$ if and only if $(D,E)$ is good in degree $k$ for every normed space $D$, if and only if $\dim_{\Kh}\widehat E<\infty$. If $K$ is complete and $k\ge p$, two
spaces with equivalent nonarchimedean norms form a pair good in degree $k$ if and only if one of them is
finite-dimensional.
\end{enumerate}
For the pairs in \textup{(1)}--\textup{(3)} and in the last sentence of \textup{(4)}, the good degrees form
$\{1,\dots,p-1\}$ or $\{1,2,\dots\}$, and every good pair has a splitting decomposition by absolute or relative
ordering splittings or by finite-coordinate splittings.
\end{introtheorem}

Theorem~\ref{intro:thm:E} collects Theorems~\ref{cls:thm:symmetric}, \ref{cls:thm:rnorm} and~\ref{fib:thm:local},
Corollaries~\ref{cls:cor:lr-table} and~\ref{cls:cor:dichotomy}, and Remark~\ref{cls:rem:mechanisms}. Items~(1)
and~(2), the last sentence for the pairs in them, and Theorem~\ref{intro:thm:D}(2) hold for spreading sequence spaces,
whose norms need only be invariant under increasing relabelings of the unit vectors (Theorem~\ref{cls:thm:spreading});
this class contains spaces with no equivalent symmetric norm (Example~\ref{cls:ex:garling}).

\begin{introtheorem}[goodness without splitting decompositions]\label{intro:thm:F}
\begin{enumerate}
\item Over every separable complete nonarchimedean field $K$ there are nonzero infinite-dimensional Banach spaces
$\fX_1,\fX_2,\dots$ with $\cL(\fX_\alpha,\fX_\alpha)=K\,\id$ and $\cL(\fX_\alpha,\fX_\beta)=0$ for $\alpha\ne\beta$.
\item Over $\kappa((X))$ with $\kappa$ finite, suitable choices in \textup{(1)} give $\fX=\fX_1$ with
$(\fX,\fX)$ good in every degree, while for every $k\ge p$ neither $(\mathrm S^k_{\fX})$ nor $(\mathrm R^k_{\fX})$
holds and $(\fX,\fX)$ has no splitting decomposition of degree $k$.
\item Over $\kappa((X))$ with $\kappa$ finite and these choices, the completed $\ell^1$-sum $E$ of the spaces $\fX_\alpha$
gives an $\ell^1$ block assembly $(E,E)$ that is good for Banach values in every degree and has no splitting
decomposition of any degree $k\ge p$, even with parameter pieces.
\end{enumerate}
\end{introtheorem}

Theorem~\ref{intro:thm:F} collects Theorems~\ref{bey:thm:rigid}, \ref{bey:thm:rigid-good} and~\ref{bey:thm:rigid-sum},
proved in Appendix~\ref{app:rigid}.

Bundles need goodness only for triples $(P;E,D)$, along analytic operator families on open subsets of the base model
$P$ (Section~\ref{sec:families}). Goodness of a triple reduces to homogeneous polynomial families and to familywise classes
(Theorems~\ref{fam2:thm:reduction} and~\ref{fam2:thm:class}). Under~\textup{(L)}, $(\ell^2(\N,K);E,E)$ with
$E=\ell^r(\N,K)$, $2<r<\infty$, is good in every degree, although neither is $\ell^2(\N,K)$ polarizable nor $(E,E)$
good for Banach values in degree $k$ (Theorems~\ref{fam2:thm:polarizable} and~\ref{fam2:thm:gamma-split},
Proposition~\ref{fam2:prop:fibers}, Corollary~\ref{cls:cor:symmetric-examples}, Theorem~\ref{fib:thm:ell2}). For a
structure group with an open subgroup that lies in a linear space of operators and contains a ball of it around the
identity, goodness is that of the linear space, not that of the closed linear span of the group
(Theorem~\ref{fam2:thm:groups}, Proposition~\ref{fam2:prop:span}).

\begin{introtheorem}[what bundles need]\label{intro:thm:G}
Let $P$ be a normed space and $k\ge1$.
\begin{enumerate}
\item Let $\cC$ be a class of Banach spaces closed under finite direct sums. Fiberwise alternating maps with Banach value
fibers preserve the analyticity of bundles with typical fibers in $\cC$ over every analytic manifold modeled on $P$ if
and only if $x\mapsto\Lambda^k_\pi u(x)$ is analytic for every analytic family $u$ of operators between spaces of $\cC$
on an open subset of $P$. For a single Banach fiber, families of invertible operators give the same condition.
\item The condition holds for all normed fibers if $P$ is polarizable in every degree $n\le k$, and over every $P$ for
fibers $E,D$ with $(E,\widehat D)$ good for Banach values in degree $k$. In characteristic $p$, the condition in all
degrees $j\le k$ is decided by the lifting constants of homogeneous polynomial families in the degrees $p^i\le k$.
Under~\textup{(L)}, it holds for all normed fibers over $P=\ell^r(\N,K)$, $1\le r<\infty$, if and only if $r<2$, and it
fails for some normed fibers over every $P$ with an equivalent nonarchimedean norm and $\dim_{\Kh}\widehat P=\infty$,
such as $c_0(\N,K)$.
\item Let $K$ be complete and $P$ a Banach space with $P\cong P^3$. The bundles $\widehat\Lambda^k_\pi TM$ are analytic
for every analytic manifold $M$ modeled on $P$ if and only if the condition holds for $E=D=P$. Under~\textup{(L)},
$c_0(\N,K)$ and $\ell^r(\N,K)$, $4<r<\infty$, model manifolds $M$ with a nonanalytic bundle $\widehat\Lambda^k_\pi TM$.
\end{enumerate}
\end{introtheorem}

Theorem~\ref{intro:thm:G} collects Proposition~\ref{fam2:prop:reductions}, Corollary~\ref{fam2:cor:bundles},
Theorems~\ref{fam2:thm:reduction}, \ref{fam2:thm:polarizable} and~\ref{fam2:thm:bases},
Proposition~\ref{fam2:prop:fibers} and Theorem~\ref{fam2:thm:tangent}.

Section~\ref{amb:sec} treats ambient analyticity as a structure, not a class of maps. Fixed envelopes come first: in every
characteristic the transitions act analytically on $\Mult^k$, preserving the alternating forms, and this gives
n\textsuperscript{o}~7.8.1--7.8.4 of~\cite{Bou} on ambient sections of filtered analytic bundles
(Theorem~\ref{amb:thm:seven-eight}). If the enlargement may depend on the map, the maps so obtained are not a usable class:
over every field of characteristic $p>0$ they are not closed under composition (Theorem~\ref{amb:thm:failure}), and the
mechanism is that enlarging the target carries maps into the space, not analytic maps out of it
(Section~\ref{amb:ssec:mechanism}). For $K$ complete and scalar values, $\Mult^k(-;K)$ is initial among the functorial
envelopes of $\Alt^k(-;K)$ with a multilinear action (Theorem~\ref{amb:thm:initial}). On tame pairs, such as ultrametric
spaces or $\ell^1(J,K)$, ambient analyticity has composition and inverse functions, and there it coincides with
divided-power analyticity, which takes bounded polynomial laws as Taylor coefficients
(Theorems~\ref{amb:thm:tame-calculus} and~\ref{amb:thm:tame-examples}). A sequel~\cite{II} studies divided-power
analyticity.

Section~\ref{sec:questions} collects questions; those on ambient analyticity are at the end of Section~\ref{amb:sec}, and
the proofs of that section are in Appendix~\ref{amb:app}. Every theorem, proposition, lemma and corollary of the paper is proved
in the text, except Theorem~\ref{thm:ingleton} and Lemma~\ref{rig:lem:ramsey}, which are cited;
Section~\ref{sec:verification} lists the statements that also have Lean proofs of corresponding statements.
Sections~\ref{sec:bundles}, \ref{ob:ssec:powers}, \ref{ob:ssec:tests}, \ref{spl:ssec:one-value},
\ref{spl:ssec:parameter}, \ref{spl:ssec:bundle}, \ref{obs:ssec:counterexamples} and~\ref{sec:geometry} with
Appendices~\ref{app:obstruction-preliminaries}--\ref{app:coordinate-methods} form the formalized core; besides it, they use
only Definition~\ref{ob:def:sequences}, the sign $\epsilon(S)$ of Section~\ref{cmb:ssec:degrees} and condition~\textup{(L)} of
Section~\ref{obs:ssec:field}. Proofs
that would interrupt the text are collected in Appendix~\ref{app:deferred}.

\subsection*{Related work}
The bundle construction follows Lang's functors of class $C^n$~\cite[Ch.~III, \S4]{La} and Bourbaki's vector
functors~\cite[n\textsuperscript{o}~7.6]{Bou}. Bourbaki defines the alternating bifunctor when $K$ has characteristic
$0$ or the bundles have finite rank, and states that it is not known whether this restriction can be
removed~\cite[n\textsuperscript{o}~7.8]{Bou}; the paragraph \emph{Bourbaki's question} in Section~\ref{sec:geometry}
shows that it cannot over complete fields of characteristic $p>0$ that are not spherically complete. Gou\"ezel's
construction~\cite{G1,G2} covers characteristic zero; the finite-dimensional case in positive characteristic, noted
in~\cite{G1}, is Proposition~\ref{prop:split-lift}(2) over complete fields. Closed-range reflection of smoothness is
formalized in~\cite{PR}. The nonarchimedean extension mechanism is a norm-controlled form of spherical
completeness~\cite{In,Schi,Sch,Yuan}. Completed nonarchimedean exterior powers go back to Serre~\cite{Se} and
Gruson~\cite{Gr}, who do not study sections of the quotient map. The rigid spaces of Section~\ref{sec:beyond} transfer a construction of Kalton and Roberts~\cite{KR}.
Corollary~\ref{ob:cor:z-linear} detects the class by the distance of a $z$-linear map from the linear maps, as for
twisted sums~\cite{Ka,KP,CSC}.

\subsection*{Conventions}
$\cL(E,F)$ denotes bounded linear maps with the operator norm, $\Mult^k(E;F)$ bounded $k$-linear maps, and
$\Alt^k(E;F)$ those vanishing whenever two arguments coincide; this condition, not antisymmetry, defines alternation
in characteristic two. The degree-zero space is $F$. Products carry the maximum norm. $E$ denotes a source fiber, $D$
a target fiber, $F$ a value space (a \emph{target fiber} in Corollary~\ref{cor:bundle-cases} and
Appendix~\ref{app:laurent}), $G$ an intermediate space, and $P$ a base model or parameter space. \emph{Finite
coordinates} means a bounded linear isomorphism with $K^d$; equivalently, a basis $e_1,\ldots,e_d$ with bounded coordinate functionals $e^*_1,\ldots,e^*_d$, which we also call finite continuous coordinates $(e_i,e^*_i)$; over a complete field this is finite dimension. Hats
denote completions (Section~\ref{sec:duality}). Exponents are written $r$; $\ell^r$, lower $r$-estimates, $r$-norms
and $\ell^2$-domination of a family are defined in Section~\ref{sec:positive}, symmetric sequence spaces in
Definition~\ref{obs:def:sequence-spaces}. A finite field $\kappa$ has $\abs\kappa$ elements. Condition~\textup{(L)} is
fixed in Section~\ref{sec:sharpness}. A norm is \emph{nonarchimedean} if $\norm{x+y}\le\max(\norm x,\norm y)$. Every
normed field of characteristic $p>0$ is nonarchimedean: Frobenius gives $\abs{x+y}^{p^m}\le2\max(\abs x,\abs y)^{p^m}$; take roots
and let $m\to\infty$.

$C^n$ ($n\le\infty$) means iterated Fr\'echet differentiability; analyticity means local expansion in bounded
multilinear coefficients with a positive norm-controlled radius. On manifolds, $C^\omega$ also requires analytic
iterated derivatives, as in Mathlib~\cite{docs}. On open sets the two agree, even for incomplete ranges: the
derivative series converges in the completion to the Fr\'echet derivative, which exists in the operator space since
nearby points are points of analyticity.
